\documentclass{article}
\usepackage{pathfinder-readable}

\title{Distributed CVaR and the Cost of Certifying Aggregate Risk}

\begin{document}
\maketitle

\begin{abstract}
The pair joins a survey of learning-augmented algorithms with a paper on distributed optimisation of
conditional value at risk (CVaR). We asked whether the distributed algorithm's output can support a costed
guarantee for the total loss faced in deployment. A fresh audit gives a conditional risk certificate, but the
algorithm's local objective can favour a decision that is worse for total risk. The thread ended as a draft
because the benchmark gap and the cost of obtaining joint observations remain unresolved for any specified
application.
\end{abstract}

\section{Introduction}

Paper Q surveys algorithms that use predictions while retaining stated performance guarantees \cite{Q}. It
asks what happens when a prediction is produced, queried, and passed to another decision rule. Its account of
composition requires a relation between the components' benchmarks, as well as a way to charge information and
other costs. A guarantee for one component does not by itself describe the cost of the whole system.

Paper P studies agents that jointly choose a decision while each sees only sampled values of its own loss
\cite{P}. Each agent measures risk by CVaR, the average loss in a specified worst tail of its distribution.
The agents exchange decisions across a changing network and use sampled loss values in place of gradients. P
proves a finite-time bound on the expected error in the average of their local CVaRs. Its projection also
keeps decisions and sampled perturbations inside the permitted decision set.

The connection was to use P's chosen decision in an operation whose loss is the sum of the agents' losses. Q's
questions then become concrete: Does P's local risk bound control risk for that sum? Can the risk of the
chosen decision be checked from new samples? What must be paid for those checks? The operation here uses the
same decision and losses as P. A different downstream policy would need its own link between decisions, costs,
and benchmarks.

\section{What was done}

The peers first separated P's two kinds of guarantee. Its finite-time theorem bounds an expectation over the
random training run. It does not certify the risk of the particular output from one run. Its projection does
give exact membership in the decision set, but membership is different from meeting a limit on tail risk. The
peers also checked P's fixed-batch last-iterate argument: it bounds the limiting true objective under its
assumptions, yet supplies no finite-time certificate for a realised output.

They next froze P's output and considered fresh validation samples at that decision. They used the
empirical-distribution bound and CVaR stability result in P's appendix to bound each local CVaR. P's
subadditivity observation then turned the local bounds into an upper bound for the CVaR of the sum, provided
all agents use the same tail level. The samples had to be independent of the training transcript: reusing the
samples that selected the decision would not justify the same conditional argument. They added a known safe
fallback and a margin assumption to estimate how often the audit would reject the proposed decision.

An objection shifted the work from estimation to information and benchmarks. Separate local samples need not
reveal how losses move together in deployment. The peers constructed examples in which this missing dependence
made a local certificate very loose, and another in which minimising the local objective chose the wrong
decision for aggregate risk. They then derived an explicit benchmark gap and considered a fresh audit of
paired losses from the actual joint law. Ada's checks are recorded in \url{../ledger.jsonl}; Emmy's derivation
and cross-check are in \url{../emmy/costed_cvar_certificate.md}.

\section{Results}

Let \(m\) be the number of agents and \(\mathcal X\) their common feasible decision set. At decision \(x\),
let \(J_i(x,\xi_i)\) be agent \(i\)'s loss, with \(\xi_i\) its random input. For a common tail probability
\(\alpha\), write \(C_i(x)=\operatorname{CVaR}_{\alpha}(J_i(x,\xi_i))\) and \(\mathcal C(x)=m^{-1}\sum_i
C_i(x)\). Let \(\mathcal C^*\) be the minimum of \(\mathcal C\) over \(\mathcal X\). P's weighted average
output \(X\) belongs to \(\mathcal X\), and its theorem supplies a deterministic bound \(B_T\) after \(T\)
iterations such that
\[
\mathbb E[\mathcal C(X)-\mathcal C^*]\le B_T.
\]
The expectation covers the random optimisation run. The bound includes the effects of network disagreement,
smoothing, and finite sampling \cite{P}.

Suppose deployment at \(x\) has total loss \(S(x,\xi)=\sum_i J_i(x,\xi_i)\), where \(\xi\) collects the
agents' inputs under a specified joint law. Write \(A(x)=\operatorname{CVaR}_{\alpha}(S(x,\xi))\), and let
\(A^*\) be its minimum over \(\mathcal
X\). CVaR subadditivity gives \(A(x)\le m\mathcal C(x)\) at every decision. Thus the nonnegative benchmark gap
\(D=m\mathcal C^*-A^*\) gives
\[
\mathbb E[A(X)-A^*]\le mB_T+D.
\]
This is the benchmark relation called for by Q's composition discussion \cite{Q}. P controls \(B_T\), while
neither P's local theorem nor a local audit bounds \(D\).

The gap can be real decision loss. Take two agents, tail probability \(\alpha=1/2\), decisions \(x\in[-1,1]\),
a fair shared indicator \(Z\), and a positive loss scale \(U\). The peers used
\[
J_1(x,Z)=UZ,\qquad
J_2(x,Z)=U\bigl((1-x)Z+x(1-Z)\bigr).
\]
The local CVaRs are \(U\) and \(U\max\{1-x,x\}\). Their sum is smallest at \(x=1/2\). The aggregate CVaR is
\(A(x)=U(2-x)\), which is smallest at \(x=1\). Choosing the exact local optimum therefore loses \(U/2\)
against the aggregate optimum, and here \(D=U/2\). Ada derived the example and Emmy independently checked it
in the files cited above.

A selective certificate is available under extra information assumptions. After freezing \(X\), draw \(n_i\)
fresh, independent validation losses for each agent at that decision. Assume their absolute values are at most
known constants \(U_i\), and let \(\widehat C_i(X)\) be the empirical CVaR. For a chosen failure probability
\(\zeta\), define
\[
\epsilon_i=\frac{2U_i}{\alpha}
\sqrt{\frac{\log(2m/\zeta)}{2n_i}},
\qquad
R=\sum_{i=1}^{m}\bigl(\widehat C_i(X)+\epsilon_i\bigr).
\]
P's CVaR stability lemma and the empirical-distribution inequality give \(A(X)\le R\) with probability at
least \(1-\zeta\) over validation data \cite{P}. Thus a rule that accepts \(X\) only when \(R\le b\), for a
fixed risk budget \(b\), falsely accepts an action with \(A(X)>b\) with probability at most \(\zeta\). On the
same validation event, a fresh deployment loss exceeds \(R\) with probability at most \(\alpha\). These are
distinct probabilities: one concerns the audit, the other a future loss draw.

The peers also checked a fallback bound. Suppose a known decision \(x_0\in\mathcal X\) satisfies \(A(x_0)\le
b\), a known positive margin \(h\) satisfies \(m\mathcal C^*\le b-h\), and the validation sizes give
\(2\sum_i\epsilon_i\le h/2\). Deploy \(X\) on acceptance and \(x_0\) otherwise. Markov's inequality applied to
P's nonnegative objective gap gives
\[
\Pr\{\text{reject }X\}\le\zeta+\frac{2mB_T}{h}.
\]
The chance that the deployed decision has aggregate CVaR above \(b\) is at most \(\zeta\), since the fallback
is known to be safe. If \(\Delta=\mathcal C(x_0)-\mathcal C^*\), its expected local-objective gap is at most
\(B_T+\Delta(\zeta+2mB_T/h)\). The margin and fallback are additional premises, not results of P.

\section{Objections and limits}

The sum of local CVaRs can overstate aggregate CVaR even with unlimited local data. In the peers' example,
exactly one of \(m\) agents loses \(U\) in each equally likely outcome. At \(\alpha=1/m\), each local CVaR is
\(U\), but the total loss is always \(U\). Dependence is also invisible to a marginal-only oracle: two fair
losses with identical local distributions can move together, giving aggregate CVaR \(2U\), or alternate,
giving aggregate CVaR \(U\). Ada and Emmy checked that this remains an information obstacle at an exact
optimum of P's local objective, provided the oracle does not supply paired outcomes. Related dependence
questions already appear in risk aggregation work \cite{cheung2010}.

A fresh audit of paired loss vectors from the actual deployment law can estimate \(A(X)\) directly. If \(n\)
such vectors are independent, their total loss is bounded in absolute value by \(U_\Sigma=\sum_i U_i\), and
\(\widehat A(X)\) is their empirical aggregate CVaR, the same argument gives an upper certificate
\[
R_J=\widehat A(X)+\frac{2U_\Sigma}{\alpha}
\sqrt{\frac{\log(2/\zeta)}{2n}}
\]
with confidence \(1-\zeta\). This removes local-sum slack at the chosen decision. It neither finds \(A^*\) nor
removes \(D\) from the decision-quality bound. Paired sensing and communication have to be available and
priced. Existing work on CVaR confidence bounds also means the holdout step alone is no novelty claim
\cite{liang2023,torossian2019}.

Finite samples cannot make this filter exactly risk-feasible for every accepted action without further
information. The peers gave a bounded, convex loss with a known safe fallback for which a safe law and an
unsafe law can produce the same finite run of zero observations. A rule accepting the candidate on that run
has a positive chance of accepting it under the unsafe law. The confidence failure term therefore cannot
simply be erased.

Finally, Q charges calls to a predictor, while P samples loss values and exchanges decisions. Charging \(N\)
training and validation evaluations at \(\kappa N\), where \(\kappa\) is a declared cost per evaluation, is a
new accounting convention \cite{Q,P}. If a query physically plays a perturbed decision, its incurred loss
matters too. Forming the shared output, collecting audit reports, and coordinating deployment may add cost. No
competitive ratio for a separate downstream policy follows without a relation for that policy's benchmark and
its sensitivity to \(X\).

\section{Why the thread stopped}

The verifier left the thread at DRAFT. It accepted the benchmark-gap argument and the example where perfect
optimisation of P's objective still loses \(U/2\) in aggregate CVaR. It also judged the audit certificate
correctly conditional and the cost and information limits acknowledged. The present evidence therefore gives a
sound connection, but no application-specific improvement over a safe baseline or complete costed downstream
theorem. The smallest restart is to choose an application with a priced source of paired observations and a
safe baseline, then prove or refute a useful bound on \(D\) and compare total cost with that baseline.

\bibliographystyle{plainurl}
\bibliography{references}

\end{document}
